\documentclass[11pt]{article}
\usepackage[utf8]{inputenc}
\usepackage[margin=1in]{geometry}
\usepackage{titlesec}
\usepackage{parskip}
\usepackage{fancyhdr}
\usepackage{amsmath, amssymb}
\usepackage{hyperref}

% Ultra-compact spacing
\titlespacing*{\section}{0pt}{6pt}{2pt}
\setlength{\parindent}{0pt}
\setlength{\parskip}{3pt}

\pagestyle{fancy}
\fancyhf{}
\lhead{\footnotesize STAT8101}
\chead{\footnotesize \href{http://www.cs.columbia.edu/~blei/teaching.html}{Invariance, Causality, and Applications}}
\rhead{\footnotesize \href{http://www.cs.columbia.edu/~blei/}{David M. Blei} (db2128)}
\cfoot{\thepage}

\title{\vspace{-2em}\textbf{Reader Report: Week 03} \\ \vspace{0.2em} \large Invariance, Causality, and Robustness \\ \footnotesize Bühlmann (2020), Peters et al. (2016), Peters, Janzing \& Schölkopf (2017)}
\author{\href{https://github.com/dhardestylewis}{Daniel Hardesty Lewis} (dl3645)}
\date{2026 Spring}

\begin{document}

\maketitle
\thispagestyle{fancy}
\vspace{-1em}

\section*{Summary}
This week's readings bridge the gap between statistical prediction and causal inference through the concept of \textbf{invariance}.
\begin{itemize}
    \item \textbf{Peters, Janzing, and Schölkopf (2017)} (Book, Ch 1, 2, 7) establish the \textbf{Principle of Independent Mechanisms}: causal mechanisms (e.g., $P(E|C)$) are autonomous and invariant to changes in other modules. While the joint distribution $P(C, E)$ may shift across environments, the conditional causal mechanism remains stable.
    \item \textbf{Peters, Bühlmann, and Meinshausen (2016)} (JRSS-B) operationalize this with \textbf{Invariant Causal Prediction (ICP)}. They define ``plausible causal predictors'' as subsets $S$ that yield an invariant conditional distribution $P(Y \mid X_S)$ across heterogeneous environments $\mathcal{E}$. The intersection of such sets, $\hat{S}(\mathcal{E})$, consistently identifies causal ancestors (Theorem 1), transforming causal discovery into a search for stability.
    \item \textbf{Bühlmann (2020)} generalizes this as a \textbf{risk minimization problem}, introducing \textbf{Anchor Regression}. This method relaxes strict invariance to find predictors robust to shift perturbations, allowing for \textbf{"diluted causality"} even when unobserved confounding prevents full causal discovery.
\end{itemize}

\section*{Critique \& Project Connection}
The shift from structure learning (finding the exact DAG) to invariant prediction (finding robust predictors) is highly relevant for applied work where the full causal graph is unknowable.
For my \href{https://github.com/dhardestylewis/thesis/blob/main/Blei-Invariance_Causality-2026Spring/Project_Proposal.d/Project_Proposal_Final_V21.pdf}{\textit{Predicting NIMBYism Causally}} final project, this framework justifies pooling data from different ``regimes'' (e.g., pre/post-crash housing markets) to find invariant predictors of opposition.
\begin{enumerate}
    \item Following Peters et al. (2016), I can treat years or market states as environments $\mathcal{E}$. If a feature (e.g., ``gentrification score'') predicts opposition $Y$ invariantly across 2007 and 2024, it is a strong candidate for a stable driver.
    \item Bühlmann's \textbf{Anchor Regression} offers a practical tool if the strict invariance of ICP is too strong for noisy text data. It helps identify ``diluted'' causal factors that remain predictive even under distribution shifts, aligning with the goal of robustly identifying NIMBYism drivers.
\end{enumerate}

\section*{Questions}
\begin{enumerate}
    \item Bühlmann (2020) suggests that ``diluted causality'' is useful when strict causal discovery fails. Does this imply we should prefer Anchor Regression over ICP for high-dimensional, noisy text data where unconfoundedness is unlikely?
    \item Peters et al. (2016) assume that the error distribution $F_\epsilon$ is identical across environments (Assumption 1). How sensitive is ICP to slight distributional drifts (e.g., concept drift in language use over decades) that do not strictly qualify as interventions?
\end{enumerate}

\newpage
\title{\vspace{-2em}\textbf{Project Proposal Update: Predicting NIMBYism Causally}}
\vspace{1em}
\hrule
\vspace{1em}

\section*{1. What problem am I solving?}
Investigating causal drivers of neighborhood opposition (NIMBYism) to housing. \textbf{Goal:} separate stable, \textbf{invariant} predictors of resistance from \textbf{spurious associations} driven by \textbf{environment-specific confounders}. \textbf{Aim:} find features predicting opposition consistently across the distinct \textbf{market regimes} of the 2018--2025 period.

\section*{2. Why is this problem important? To whom?}
Critical for urban planners. Current models overfit to specific eras, failing when market regimes shift. Identifying \textit{invariant} causes enables policy design addressing fundamental concerns rather than symptoms of a moment, ensuring \textbf{robustness}.

\section*{3. What is my strategy for solving it?}
Using \textbf{Invariant Risk Minimization (IRM)} and \textbf{Bayesian Invariant Prediction (BIP)}, grounded in the \textbf{Invariant Prediction} framework:
\begin{enumerate}
    \setlength\itemsep{-0.3em}
    \item \textbf{Environments:} Partition data via \textbf{expanding window Cross-Validation} into \textbf{Training Environments} ($\mathcal{E}_{tr}$) and \textbf{Test Environments} ($\mathcal{E}_{te}$).
    \item \textbf{Representation:} Extract features $\Phi(X)$ from property data and protest letters.
    \item \textbf{Methods:} Implement two core paradigms:
    \begin{itemize}
        \item \textbf{Invariant Risk Minimization (IRM):} A \textit{regularized optimization} paradigm to learn $\Phi$ such that a classifier is simultaneously optimal across all $\mathcal{E}_{tr}$.
        \item \textbf{Bayesian Invariant Prediction (BIP):} A \textit{probabilistic inference} paradigm defining a latent feature selection vector $z$; the posterior $p(z|D)$ targets the true invariant features.
    \end{itemize}
\end{enumerate}

\section*{4. Why is this solution good? Where does it fall short?}
\textbf{Falls short:} It assumes an \textbf{invariant causal mechanism}. The assumption fails if the underlying reasons for NIMBYism undergo \textbf{mechanism change} over time (e.g., shifting from traffic concerns to displacement concerns).

\section*{5. How can I demonstrate that the solution works?}
Train on expanding $\mathcal{E}_{tr}$, test on all \textbf{future horizons} in $\mathcal{E}_{te}$. A successful Invariant Model demonstrates \textbf{predictive stability} through the \textbf{invariance of its AUC} score across all horizons, maintaining reliability where baselines fail under \textbf{distribution shift}.

\end{document}
